% !TeX root = ../../../main.tex
\FloatBarrier
\subsection{Paired matched-mention sensitivity analysis}\label{ch08:supp:matched-mention}
\noindent
\textbf{Methods.} This secondary analysis paired the same manually annotated reference mention across unimodal and multimodal reconstruction. Event recovery was evaluated over all 8,066 reference mentions with exact two-sided McNemar tests. For temporal accuracy, we retained only reference events matched under both conditions, compared paired absolute timestamp errors with two-sided Wilcoxon signed-rank tests, and recomputed AULTC on the same shared matched subset. Wilcoxon calculations used tie correction and dropped zero differences. No multiplicity adjustment was applied. Because mentions are clustered within 40 discharge summaries and event-rich cases contribute more observations, the reported p values are nominal and may be anti-conservative. These analyses are exploratory and secondary to the paired case-level bootstrap in \hyperref[ch08:tab:supp-paired-bootstrap]{Supplementary Table~\ref*{ch08:tab:supp-paired-bootstrap}}.

\noindent
\textbf{Results.} Event-recovery differences were small: DeepSeek R1 and DeepSeek V3.2 were unchanged, GPT-OSS-120B changed negligibly, GLM5 and Qwen3.5-397B-A17B increased slightly, and KimiK2-Instruct decreased slightly. On the shared matched subset, AULTC point estimates favored multimodal calibration for DeepSeek V3.2, GPT-OSS-120B, GLM5, and Qwen3.5-397B-A17B; DeepSeek R1 changed negligibly and KimiK2-Instruct decreased. The nominal mention-level p values should not be interpreted as confirmatory evidence because they do not model within-case dependence.

\par\medskip
\noindent
% !TeX root = ../../../main.tex
\begin{table}[htbp]
\caption{Paired matched-mention event-recovery sensitivity analysis at event-matching threshold 0.1. Values use the same pooled reference mentions under both conditions, with Delta = multimodal minus unimodal.}\ChEightFloatLabel{ch08:tab:supp-mcnemar}
\ChEightTableBodyStart
\begin{ChEightFloatBody}

\small

\centering
\begin{tabular}{@{}lcccc@{}}
\toprule
\textbf{Model} & \textbf{Unimodal} & \textbf{Multimodal} & \textbf{Delta} & \textbf{Nominal McNemar p} \\
\midrule
DeepSeek R1 & 0.39 & 0.39 & 0.00 & 1.000 \\
KimiK2-Instruct & 0.26 & 0.25 & -0.01 & $<0.001$ \\
GPT-OSS-120B & 0.32 & 0.32 & 0.00 & 0.558 \\
DeepSeek V3.2 & 0.49 & 0.49 & 0.00 & 1.000 \\
GLM5 & 0.33 & 0.33 & 0.00 & $<0.001$ \\
Qwen3.5-397B-A17B & 0.38 & 0.40 & 0.02 & $<0.001$ \\
\bottomrule
\end{tabular}
\par
\begin{minipage}{\textwidth}
\footnotesize\textit{Note:} Exact two-sided McNemar tests compare paired match status over all 8,066 reference mentions. Values are rounded to two decimals, so a displayed Delta of 0.00 may reflect a small nonzero unrounded difference. The nominal p values do not account for clustering of mentions within discharge summaries and are reported only as an exploratory sensitivity analysis.
\end{minipage}

\end{ChEightFloatBody}
\end{table}
\par\medskip

\par\medskip
\noindent
% !TeX root = ../../../main.tex
\begin{table}[htbp]
\caption{Paired temporal-accuracy sensitivity analysis on reference events matched under both unimodal and multimodal reconstruction. Delta AULTC = multimodal minus unimodal.}\ChEightFloatLabel{ch08:tab:supp-wilcoxon}
\ChEightTableBodyStart
\begin{ChEightFloatBody}

\small

\centering
\begin{adjustbox}{max width=\textwidth}
\begin{tabular}{@{}lccccc@{}}
\toprule
\textbf{Model} & \textbf{Shared matches, n} & \textbf{AULTC (UM)} & \textbf{AULTC (MM)} & \textbf{Delta AULTC} & \textbf{Nominal Wilcoxon p} \\
\midrule
DeepSeek R1 & 3,134 & 0.80 & 0.79 & 0.00 & 0.364 \\
KimiK2-Instruct & 1,999 & 0.68 & 0.68 & -0.01 & $<0.001$ \\
GPT-OSS-120B & 2,532 & 0.76 & 0.77 & 0.01 & $<0.001$ \\
DeepSeek V3.2 & 3,930 & 0.71 & 0.77 & 0.06 & $<0.001$ \\
GLM5 & 2,650 & 0.78 & 0.78 & 0.00 & 0.003 \\
Qwen3.5-397B-A17B & 3,025 & 0.74 & 0.77 & 0.03 & $<0.001$ \\
\bottomrule
\end{tabular}
\end{adjustbox}
\par
\begin{minipage}{\textwidth}
\footnotesize\textit{Note:} Nominal p values are from paired two-sided Wilcoxon signed-rank tests of absolute timestamp error among events matched in both conditions; AULTC is reported on the same shared subset. The tests do not account for clustering within discharge summaries and may be anti-conservative. UM = unimodal; MM = multimodal.
\end{minipage}

\end{ChEightFloatBody}
\end{table}
\par\medskip

